\section{Integración del sistema y flujo de despliegue}
\subsection{Slide 09-10 Pipeline}
Las Slides 09-10 sintetizan el sistema Autonomy: Reward module, Reset policy, Goal manager y Safety guard trabajan por capas, y los resultados se envían al dashboard para que el operator los revise.

\begin{figure}[H]
\centering
\includegraphics[width=0.9\textwidth]{slides-images/slide-10.jpg}
\caption{La Slide 10 muestra la modular pipeline del sistema Autonomy.}
\end{figure}

\begin{importantbox}{Nodos clave del pipeline de despliegue}
\begin{itemize}
\item Cada Reward model update desencadena un shadow policy test y se hace un sanity check en simulation;\\
\item El Reset controller ejecuta todos los días una recovery routine para asegurar que la forward-backward policy siga siendo confiable;\\
\item El Goal manager actualiza la goal library con base en las live metric y descarta todo goal stale o unsafe;\\
\end{itemize}
\end{importantbox}

\subsection{Retroalimentación entre equipos e intercambio de conocimiento}
El ponente advierte: Autonomy necesita ciclos de retroalimentación entre equipos; por ejemplo, el operator sube cada reset failure case al shared log y, a partir de ello, el equipo de policy ajusta la goal selection o la reward normalization. La documentación y la weekly review meeting ayudan a mantener el conocimiento.

\subsection{Resumen de la sección}
La integración del sistema se centra en organizar el pipeline y en la comunicación entre equipos, para asegurar que Autonomy no solo sea eficaz en la investigación, sino que también responda con rapidez a las anomalías durante el despliegue.
\newpage
